\section{תוצאות: חיזוי יום־קדימה}

\subsection{תמונה כללית}

טבלה \ref{tab:meanDPA} מציגה את ממוצע ה־DPA הבלתי־משוקלל בין תשעת הנכסים. התוצאה החשובה ביותר היא שאין מודל אחד שמציג יתרון חיזויי עקבי.

\begin{table}[H]
\centering
\begin{tabular}{l r}
\toprule
Model & Mean DPA (\%) \\
\midrule
Naive - train mean & 54.54 \\
Logistic Regression & 54.14 \\
Random Forest Classifier & 53.27 \\
Gaussian HMM - soft posterior & 53.06 \\
HistGradientBoostingClassifier & 52.49 \\
Gaussian HMM - hard state & 52.06 \\
\bottomrule
\end{tabular}
\caption{ממוצע DPA בין תשעת הנכסים. הממוצע אינו שקול למבחן מובהקות ואינו מוכיח עליונות כללית של אף מודל.}
\label{tab:meanDPA}
\end{table}

\begin{figure}[H]
\centering
\includegraphics[width=0.98\textwidth]{dpa_comparison_all_assets.png}
\caption{השוואת DPA בין Baseline, שלושת מודלי ה־Supervised ML ו־HMM hard/soft לכל נכס.}
\label{fig:dpaall}
\end{figure}

\subsection{תוצאות לפי נכס}

יש כמה מקרים מקומיים מעניינים. ב־\asset{QQQ}, Random Forest השיג \pct{59.22}, soft HMM השיג \pct{58.98}, וה־Naive baseline \pct{58.74}. ב־\asset{BTC-USD}, Random Forest השיג \pct{52.94} וה־soft HMM \pct{52.05}, לעומת \pct{51.34} ל־Naive. לעומת זאת, ב־\asset{JPM} ה־Naive הגיע ל־\pct{57.52} בעוד soft HMM הגיע ל־\pct{51.70}. ב־\asset{TLT}, HistGradientBoosting היה הטוב ביותר מבין המודלים שנבדקו עם \pct{53.40}, בעוד soft HMM הגיע ל־\pct{47.09}.

הפיזור הזה חשוב יותר מהדירוג הממוצע: HMM אינו נכשל בכל נכס, אך גם אינו מנצח בצורה יציבה. אותו דבר נכון למודלי ה־Supervised ML. לכן אין evidence לכך שה־Feature Vector הנוכחי מכיל signal חזק ועקבי לחיזוי next-day direction.

\subsection{למה DPA לבדו עלול להטעות}

ב־\asset{SPY}, למשל, Naive train-mean ו־HMM מגיעים ל־DPA של כ־\pct{58.50}. עם זאת, Logistic Regression מגיע ל־DPA של \pct{58.01} אך Balanced Accuracy של כ־\pct{49.6} ו־ROC-AUC של כ־\pct{47.3}. Random Forest משיג ROC-AUC של כ־\pct{50.4}. כלומר חלק ניכר מה־Accuracy נובע מחוסר איזון בין ימי Up ו־Down, ולא מיכולת חזקה לדרג ימים לפי הסתברות לעלייה.

זו הסיבה שבדוח לא מתייחסים ל־DPA מעל \pct{50} כהוכחה אוטומטית ל־predictive edge. מודל שמנבא את המחלקה הנפוצה יכול לקבל Accuracy סביר גם בלי להבין את דינמיקת היום הבא.

\subsection{Hard state לעומת Soft posterior}

בממוצע, soft HMM משיג \pct{53.06} לעומת \pct{52.06} ל־hard-state. הפער קטן, אך הוא עקבי עם האינטואיציה שלפיה שימוש בכל וקטור ה־posterior שומר מידע שנזרק ב־argmax. לדוגמה, ב־\asset{QQQ} hard HMM הגיע ל־\pct{56.31} בעוד soft HMM הגיע ל־\pct{58.98}; ב־\asset{BTC-USD} הפער גדול יותר, \pct{48.13} לעומת \pct{52.05}. מנגד קיימים גם נכסים שבהם soft אינו משפר.

לכן המסקנה אינה ש־soft posterior ``פותר'' את forecasting, אלא שהוא דרך הגיונית יותר לנצל posterior uncertainty מאשר hard assignment כאשר משתמשים ב־HMM לתחזית.

\subsection{Walk-forward robustness על SPY}

שלושת ה־walk-forward folds מציגים תמונה מעורבת. soft HMM השיג בהתאמה \pct{50.36}, \pct{51.82} ו־\pct{60.22}; Naive train-mean השיג \pct{48.54}, \pct{52.92} ו־\pct{59.85}. הממוצע הוא \pct{54.14} ל־soft HMM לעומת \pct{53.77} ל־Naive, אך HMM אינו מנצח בכל fold והפער הכולל קטן.

\begin{figure}[H]
\centering
\includegraphics[width=0.82\textwidth]{walk_forward_dpa.png}
\caption{DPA בשלושה expanding walk-forward folds על SPY. ה־HMM אינו מפגין יתרון עקבי בין התקופות.}
\label{fig:walkforward}
\end{figure}

בדיקת ה־walk-forward מחזקת את הצורך בניסוח זהיר: ייתכן שמידע על regime מסייע בתקופות מסוימות, אך לא נמצא predictive advantage יציב שניתן להכליל לכל תקופה.
